\begin{name}
	{\tenchude}
	{TOÁN 10}
	{LỚP TOÁN THẦY PHÁT}
	{Thời gian: 90 phút - Không kể thời gian phát đề}
\end{name}
\Opensolutionfile{ansbook}[ans/ansbookDe4]
\TN
\Opensolutionfile{ans}[ans/ansDe4-TN1]
\begin{ex}%[Mức 2]Giảng Toán 10, Nguyễn Tài Tuệ]%[0D3H1-5]
	Tổng tất cả các giá trị nguyên dương của tham số $m$ để hàm số $y=-2 x^2+(m+1) x+3$ nghịch biến trên khoảng $(1 ; 5)$ là
	\choice
	{\True $6$}
	{$3$}
	{$1$}
	{$15$}
	\loigiai
	{
		Ta có hoành độ đỉnh $ \dfrac{-b}{2a}=\dfrac{m+1}{4} $.\\
		Để hàm số nghịch biến trên khoảng $ (1;5) $ thì $ \dfrac{m+1}{4} \le 1 \Leftrightarrow m \le 3 $.\\
		Suy ra các giá trị nguyên dương của $m$ thỏa là $ 1 $, $ 2 $, $ 3 $.\\
		Tổng các giá trị của $ m $ là $ 1+2+3=6 $.
	}
\end{ex}

\begin{ex}%[24-25 giảng K10 - K11, Phan Quốc Trí]%[0D3N2-1]
	Hoành độ đỉnh của parabol $(P) \colon y=2x^2-4x+3$ bằng
	\choice
	{$-2$}
	{$2$}
	{$-1$}
	{\True $1$}
	\loigiai{
		Hoành độ đỉnh của $(P)$ là $x_I=-\dfrac{b}{2a}=\dfrac{4}{2 \cdot 2}=1$.}
\end{ex}

\begin{ex}%[0D3H2-2]%[Yến Trần]%[Đề cương Nguyễn Thượng Hiền_HK1]
	Cho hàm số $y=f(x)=x^2-2x+3$. Trong các khẳng định sau, tìm khẳng định đúng.
	\choice
	{\True Hàm số nghịch biến trên khoảng $(-\infty;1)$}
	{Hàm số đồng biến trên khoảng $(0;+\infty)$}
	{Đồ thị hàm số đi qua điểm $A(1;0)$}
	{Hàm số đồng biến trên khoảng $(-1;+\infty)$}
	\loigiai
	{
		Ta có $a=1$, $b=-2$, $c=3$ nên $-\dfrac{b}{2a}=-\dfrac{-2}{2 \cdot 1}=1$.\\
		Do hệ số $a>0$ nên ta có bảng biến thiên
		\begin{center}
			\begin{tikzpicture}[>=stealth]
				\tkzTabInit[nocadre=false,lgt=1,espcl=2,deltacl=0.5]{$x$/.7 ,$y$/2}
				{$-\infty$ , $1$ , $+\infty$}
				\tkzTabVar{+/$+\infty$ , -/$2$ , +/$+\infty$}
			\end{tikzpicture}
		\end{center}
		Suy ra hàm số nghịch biến trên khoảng $(-\infty;1)$.
	}
\end{ex}

\begin{ex}%[10-11-GK2-2425]%[Đỗ Minh Phúc]%[0D6H1-1]
	Nếu lấy $3{,}1$ làm giá trị gần đúng của $3{,}14$ thì sai số tuyệt đối là
	\choice
	{$-6{,}24$}
	{$-0{,}04$}
	{$6{,}24$}
	{\True $0{,}04$}
	\loigiai{
	Sai số tuyệt đối là $\Delta_a=\left|3{,}14-3{,}1\right|=0{,}04$.
	}
\end{ex}

\begin{ex}%[HK1, THPT Chuyên Hùng Vương - Phú Thọ, 2023-2024]%[Lê Quốc Hiệp, 10-11EX-HK1-2324]%[0D6H4-2]
	Mẫu số liệu sau cho biết chiều cao (đơn vị cm) của các bạn trong tổ
	\[\begin{array}{llllllll}163 & 159 & 172 & 167 & 165 & 168 & 170 & 161\end{array}\]
	Khoảng biến thiên của mẫu số liệu này bằng
	\choice
	{$11$}
	{$10$}
	{\True $13$}
	{$12$}
	\loigiai
	{
		Sắp xếp mẫu số liệu theo thứ tự không giảm
		\[
			\begin{array}{llllllll}
				159 & 161 & 163 & 165 & 167 & 168 & 170 & 172
			\end{array}
		\]
		Khoảng biến thiên $R=x_8-x_1=172-159=13$.
	}
\end{ex}

\begin{ex}%[soạn để GHK, HK]%[Huỳnh Xuân Tín]%[0D0N1-2]
	Gieo hai đồng tiền một lần. Kí hiệu $S$, $N$ lần lượt để chỉ đồng tiền xuất hiện mặt sấp, đồng tiền xuất hiện mặt ngửa. Mô tả không gian mẫu nào dưới đây là đúng?
	\choice
	{$\Omega=\{S; N\}$}
	{$\Omega=\{NN; SS\}$}
	{$\Omega=\{SN; NS\}$}
	{\True $\Omega=\{SN; NS; SS; NN\}$}
	\loigiai{
		Không gian mẫu $\Omega=\{SN; NS; SS; NN\}$.
	}
\end{ex}

\begin{ex}%[0D0N1-3]%[tex hóa đề CK2 - 2025 - Huyên Nguyễn]
	Tâm đi từ nhà mình đến nhà Huyền, cùng Huyền đi đến nhà Linh chơi. Biết từ nhà Tâm đến nhà Huyền có $5$ con đường đi. Từ nhà Huyền đến nhà Linh có $7$ con đường đi. Hỏi có bao nhiêu cách để Tâm đi đến nhà Linh mà phải đi qua nhà Huyền?
	\choice{$12$}{\True$35$}{$20$}{$25$}
	\loigiai{
		\begin{itemize}
			\item Từ nhà Tâm đến nhà Huyền có $5$ cách.
			\item Từ nhà Huyền đến nhà Linh có $7$ cách.
		\end{itemize}
		Vậy theo quy tắc nhân, Tâm đi đến nhà Linh mà phải đi qua nhà Huyền sẽ có $5\cdot 7=35 $ cách.
	}
\end{ex}

\begin{ex}%[0H9N3-1]
	Vectơ nào dưới đây là một vectơ chỉ phương của đường thẳng song song với trục $Ox$?
	\choice
	{\True $\vec u=\left( 1;0 \right)$}
	{$\vec u=\left((1;-1\right)$}
	{$\vec u=\left(1;1\right)$}
	{$\vec u=\left(0;1\right)$}
	\loigiai{
		Véc-tơ $\vec i=\left(1;0\right)$ là một vectơ chỉ phương của trục $Ox$.\\
		Các đường thẳng song song với trục $Ox$ có một vectơ chỉ phương là $\vec u=\vec i=\left(1;0\right)$.
	}
\end{ex}

\begin{ex}%[10-11-12EX-GHK2-2425]%[Dương Quang]%[0H9N3-4]
	Trong mặt phẳng $Oxy$, cho hai đường thẳng $\Delta_1\colon a_1x+b_1y+c_1=0$ ($a_1^2+b_1^2>0$) và $\Delta_2\colon a_2x+b_2y+c_2=0$ ($a_2^2+b_2^2>0$). Trong các công thức sau đây, công thức nào đúng?
	\choice
	{$\cos\left(\Delta_1,\Delta_2\right)=\dfrac{\left|a_1b_1+a_2b_2\right|}{\sqrt{a_1^2+a_2^2}\cdot\sqrt{b_1^2+b_2^2}}$}
	{$\cos\left(\Delta_1,\Delta_2\right)=\dfrac{\left|a_1a_2+b_1b_2\right|}{\sqrt{a_1^2+a_2^2}\cdot\sqrt{b_1^2+b_2^2}}$}
	{\True $\cos\left(\Delta_1,\Delta_2\right)=\dfrac{\left|a_1a_2+b_1b_2\right|}{\sqrt{a_1^2+b_1^2}\cdot\sqrt{a_2^2+b_2^2}}$}
	{$\cos\left(\Delta_1,\Delta_2\right)=\dfrac{\left|a_1a_2+b_1b_2\right|}{\sqrt{a_1^2+b_1^2}\cdot\sqrt{a_2^2+b_2^2}}$}
	\loigiai{
		Ta có $\cos\left(\Delta_1,\Delta_2\right)=\dfrac{\left|a_1a_2+b_1b_2\right|}{\sqrt{a_1^2+b_1^2}\cdot\sqrt{a_2^2+b_2^2}}$.
	}
\end{ex}

\begin{ex}%[0H9N3-3]
	Trong mặt phẳng tọa độ $Oxy$, cho hai đường thẳng $d_1\colon 2x+3y+4=0$ và $d_2\colon 4x+6y-1=0$. Hãy chọn khẳng định đúng.
	\choice
	{$d_1$ và $d_2$ cắt nhau}
	{\True $d_1$ và $d_2$ song song với nhau}
	{$d_1$ và $d_2$ vuông góc nhau}
	{$d_1$ và $d_2$ trùng nhau}
	\loigiai{
		Đường thẳng $d_1\colon 2x+3y+4=0$ có vectơ pháp tuyến $\overrightarrow{n}_1=(2;3)$.\\
		Đường thẳng $d_2\colon 4x+6y-1=0$ có vectơ pháp tuyến $\overrightarrow{n}_2=(4;6)$.\\
		Ta thấy $\overrightarrow{n}_2=2\overrightarrow{n}_2$ nên hai đường thẳng song song hoặc trùng nhau.\\
		Mặt khác, ta lại có điểm $M(-2;0)$ thuộc đường thẳng $d_1$ nhưng không thuộc đường thẳng $d_2$ nên hai đường thẳng này không trùng nhau.\\
		Vậy hai đường thẳng $d_1$ và $d_2$ song song với nhau.
	}
\end{ex}

\begin{ex}%[0H9N5-1]%[Dự án đề kiểm tra Toán 10, 11 HKII NH23-24-BCTuan]%[THPT LTV-Hà Nội]
	Trong mặt phẳng tọa độ $Oxy$, cho elip có phương trình  $\dfrac{x^2}{4}+\dfrac{y^2}{1}=1$. Một tiêu điểm của elip có tọa độ là
	\choice
	{$D\left(0;\sqrt{5}\right)$}
	{$B\left(0;\sqrt{3} \right)$}
	{$C\left(\sqrt{5};0\right)$}
	{\True $A\left(\sqrt{3};0\right)$}
	\loigiai{
		Các tiêu điểm của elip có tọa độ là $(c;0)$ và $(-c;0)$.\\
		Ta có $a^2=4$, $b^2=1$ nên $c^2=a^2-b^2=3\Leftrightarrow c=\sqrt{3}$.\\
		Vậy một tiêu điểm của elip có tọa độ là $A\left(\sqrt{3};0\right)$.
	}
\end{ex}

\begin{ex}%[0H9N5-5]%[Dự án đề kiểm tra Toán 10 HKII NH23-24- Đoàn Minh Tâm]%[THPT Mạc Đĩnh Chi]
	Trong mặt phẳng tọa độ $Oxy$, hypebol $(H)$ có độ dài trục thực bằng $ 6$ và tiêu cự bằng $ 10$ có phương trình chính tắc là
	\choice
	{\True $(H)\colon\dfrac{x^2}{9}-\dfrac{y^2}{16}=1$}
	{$(H)\colon\dfrac{x^2}{16}-\dfrac{y^2}{9}=1$}
	{$(H)\colon\dfrac{x^2}{25}-\dfrac{y^2}{9}=1$}
	{$(H)\colon\dfrac{x^2}{36}-\dfrac{y^2}{64}=1$}
	\loigiai{
		Độ dài trục thực bằng $6$ nên $2a=6 \Rightarrow a=3\Rightarrow a^2=9$.\\
		Tiêu cự bằng $10$ nên $2c=10 \Rightarrow c=5$.\\
		Ta có $b^2=c^2-a^2=25-9=16$.\\
		Vậy phương trình chính tắc của hypebol $(H)$ là $\dfrac{x^2}{9}-\dfrac{y^2}{16}=1$.
	}
\end{ex}
\Closesolutionfile{ans}

\TNTF
\Opensolutionfile{ans}[ans/ansDe4-TN2]
\begin{ex}%[0D6H4-2]
	Điểm trung bình các môn trong kỳ thi tốt nghiệp trung học phổ thông năm 2024 được thống kê trong bảng sau:
	\begin{center}
		\begin{tabular}{|c|c|c|c|c|c|c|c|c|c|}
			\hline
			\textbf{Môn học} & Toán & Văn  & Vật lý & Hóa học & Sinh học & Lịch sừ & Địa lý & GDCD & Ngoại ngữ \\
			\hline
			\textbf{Điểm}    & 6,45 & 7,23 & 6,67   & 6,68    & 6,28     & 6,57    & 7,19   & 8,16 & 5,51      \\
			\hline
		\end{tabular}
	\end{center}
	Các khẳng định sau đúng hay sai?
	\choiceTF[t]
	{\True Điểm trung bình của $9$ môn thi tốt nghiệp năm $2024$ (làm tròn đến hàng phần trăm) là $6{,}75$}
	{\True Khoảng biến thiên của mẫu số liệu trên là $2{,}65$}
	{\True Quy tròn số trung bình với độ chính xác $d=0{,}01$ bằng $a=6{,}8$}
	{Sai số tương đối của số gần đúng trên xấp xỉ bằng $1,48\%$}
	\loigiai{
		\begin{itemchoice}
			\itemch Vì $\bar{x}=\dfrac{6,45+7,23+6,67+6,68+6,28+6,57+7,19+8,16+5,51}{9}=\dfrac{3037}{450} \approx 6,75$
			\itemch Giá trị bé nhất của mẫu số liệu trên là: $5{,}51$. \\
			Giá trị lớn nhất của mẫu số liệu trên là: $8{,}16$. \\
			Do đó: khoảng biến thiên của mẫu số liệu trên là $8{,}16-5{,}51=2{,}65$.
			\itemch Quy tròn số trung bình với độ chính xác $d=0{,}01$ bằng $a=6{,}8$.
			\itemch Sai số tương đối của số gần đúng
			$$\delta=\dfrac{\left|6{,}8-6{,}75\right|}{6{,}8}\approx 0{,}74\%$$
		\end{itemchoice}
	}
\end{ex}

\begin{ex}%[0H9H3-2]
	Cho ba điểm $M(1;-2)$, $N(4;-1)$ và $P(1;2)$. 
	\choiceTF
	{\True Đường thẳng $\Delta$ đi qua $M$ và $N$ có vectơ pháp tuyến là $\vec{n}=(1;-3)$}
	{\True Phương trình tổng quát của đường thẳng $\Delta \colon x-3y-7=0$}
	{Đường tròn ngoại tiếp tam giác $MNP$ có tâm là $I(1;0)$}
	{\True Phương trình đường tròn ngoại tiếp tam giác $MNP$ là $(x-2)^2+y^2=5$}
	\loigiai{
		\begin{itemchoice}
			\itemch Đường thẳng $\Delta \colon \heva{&x=1+3t\\&y=-2+t}$ có vectơ chỉ phương là $\vec{u}=\left(3;1\right)$ nên nhận $\vec{n}=\left(1;-3\right)$ là một vectơ pháp tuyến.
			\itemch Đường thẳng $\Delta$ đi qua điểm $M\left(1;-2\right)$ nên phương trình tổng quát của $\Delta$ là \[\left(x-1\right)-3\left(y+2\right)=0\Leftrightarrow x-3y-7=0.\]
			\itemch Đường tròn ngoại tiếp tam giác $MNP$ có tâm $I(x;y)$ nên 
			\[\heva{&IM=IN \\ &IM=IP}\Leftrightarrow \heva{&(x-1)^2+(y+2)^2=(x-4)^2+(y+1)^2 \\ &(x-1)^2+(y+2)^2=(x-1)^2+(y-2)^2}\Leftrightarrow \heva{&x^2+y^2-2x+4y+5=x^2+y^2-8x+2y+17 \\ &x^2+y^2-2x+4y+5=x^2+y^2-2x-4y+5}\Leftrightarrow \heva{&6x+2y=12 \\ &8y=0}\Leftrightarrow \heva{&x=2 \\ &y=0}.\]
			\itemch Đường tròn ngoại tiếp tam giác $MNP$ có tâm $I(2;0)$ và bán kính $R=IM=\sqrt{(2-1)^2+(0+2)^2}=\sqrt{5}$ nên phương trình đường tròn ngoại tiếp tam giác $MNP$ là $(x-2)^2+y^2=5$.
		\end{itemchoice}
	}
\end{ex}
\Closesolutionfile{ans}

\TNSA
\Opensolutionfile{ans}[ans/ansDe4-TN3]
\begin{ex}%[Dự án Form mới 12/4/6-HK1 Lớp 10-11]%[VN-MT-6, Võ NGuyên Thạch]%[THPT TÂN TÚC-TP HCM]%[0D3H2-3]
	Biết hàm số $y=x^2-2x-1$ nghịch biến trên khoảng $(-\infty;a)$ và đồng biến trên khoảng $(a;+\infty)$. Tìm $a$.
	\shortans[]{$1$}
	\loigiai{
		\begin{itemize}
			\begin{multicols}{2}
				\item Hàm số $y=x^2-2x-1$.
				\item Tập xác định: $\mathscr{D}=\mathbb{R}$.
				\item Trục đối xứng: $x=1$.
				\item Đỉnh $S(1;-2)$.
			\end{multicols}
			\item Bảng biến thiên.
			      \begin{center}
				      \begin{tikzpicture}[font=\normalsize,t style/.style={style=solid}]
					      %dòng khai báo
					      \tkzTabInit[nocadre=true,lgt=1,espcl=4,deltacl=0.5]
					      {$x$ /0.7, $y$/2}
					      {$ -\infty$,$1$,$+\infty$}
					      %dòng biến thiên
					      \path ($(N11)!0.2!(N12)$) node (A1){$+\infty$}
					      ($(N21)!0.9!(N22)$) node (A2){$ -2 $}
					      ($(N31)!0.2!(N32)$) node (A3){$ +\infty $};
					      \foreach \x/\y in {A1/A2,A2/A3}{
							      \draw[-stealth] (\x)--(\y);
						      }
				      \end{tikzpicture}
			      \end{center}
			\item Nhận xét: Hàm số đồng biến trên $(1;+\infty)$, hàm số nghịch biến trên $(-\infty;1)$.
		\end{itemize}
		Vậy $a=1$.
	}
\end{ex}

\begin{ex}%[0D7V2-7]
	Một cửa hàng buôn giày nhập một đôi với giá là $40$ USD. Cửa hàng ước tính rằng nếu đôi giày được bán với giá $x$ USD thì mỗi tháng khách hàng sẽ mua $\left(120-x\right)$ đôi. Hỏi cửa hàng bán một đôi giày giá bao nhiêu thì thu được nhiều lãi nhất?\par
	\shortans{$80$}
	\loigiai{
		Gọi $y$ (USD) là số tiền lãi của cửa hàng bán giày.\\
		Ta có
		\allowdisplaybreaks
		\begin{eqnarray*}
			y&=&\left(120-x\right)\left(x-40\right)\\
			&=&-x^2+160x-4800\\
			&=&-\left(x-80\right)^2+1600\le 1600.
		\end{eqnarray*}
		Dấu \lq\lq =\rq\rq\ xảy ra $\Leftrightarrow x=80$.\\
		Vậy cửa hàng lãi nhiều nhất khi bán đôi giày với giá $80$ USD.
	}
\end{ex}

\begin{ex}%[0D6H3-2]%[tex hóa đề ck2 - form 2025 - đợt 2- Duong Xuan Loi]
	Điểm thi cuối học kì II tám môn Toán, Văn, Anh, Sinh, Sử, Địa, Lý, Hóa của một học sinh lần lượt là $8$; $7{,}5$; $8{,}5$; $7$; $9$; $8$; $6{,}5$; $9{,}5$. Điểm thi trung bình tám môn thi của học sinh là bao nhiêu?
	\shortans[0]{$8,0$}
	\loigiai{
	Điểm thi trung bình tám môn thi của học sinh là $\dfrac{8+7{,}5+8{,}5+7+9+8+6{,}5+9{,}5}{8}=8{,}0$.}
\end{ex}

\begin{ex}%[0H9V4-2]
	Trong mặt phẳng tọa độ $Oxy$, cho điểm $A(-1;1)$, $B(0;-2)$, $C(0;2)$. Đường tròn ngoại tiếp tam giác $ ABC $ có phương trình là $ (x-a)^2+(y-b)^2=R^2 $. Tính $ S=a+b+R^2 $.
	\shortans{$6$}
	\loigiai{
		Giả sử tâm đường tròn là điểm $I(a;b)$.\\
		Ta có \allowdisplaybreaks
		\begin{eqnarray*}
			& &IA=IB=IC\Leftrightarrow IA^2=IB^2=IC^2\\
			&\Leftrightarrow&\heva{& (-1-a)^2+(1-b)^2=(0-a)^2+(-2-b)^2\\ & (0-a)^2+(-2-b)^2=(0-a)^2+(2-b)^2}\\
			&\Leftrightarrow& \heva{& a^2+b^2+2a-2b+2=a^2+b^2+4b+4\\ & a^2+b^2+4b+4=a^2+b^2-4b+4}\\
			&\Leftrightarrow& \heva{& 2a-2b=4b+2\\ & b=0}\Leftrightarrow\heva{& a=1\\ & b=0.}
		\end{eqnarray*}
		Đường tròn tâm $I(1;0)$, bán kính $R=IC=\sqrt{a^2+b^2-4b+4}=\sqrt{5}$.\\
		Suy ra phương trình đường tròn đi qua ba điểm $A$, $B$, $C$ có dạng $(x-1)^2+y^2=5$.\\
		Vậy $a=1$, $b=0$, $R^2=5$, $a+b+R^2=6$.
	}
\end{ex}

\Closesolutionfile{ans}

\TL
\begin{ex}%[0D3H1-2]%[Yến Trần]%[Đề cương Nguyễn Thượng Hiền_HK1]
	Tìm tập xác định của hàm số $y=\dfrac{2x+2022}{x^2-4}+\dfrac{2023x}{\sqrt{x+1}}$.
	\loigiai{
		Điều kiện xác định của hàm số đã cho là $\heva{&x^2-4\ne0\\&x+1>0} \Leftrightarrow \heva{&x\ne \pm 2\\&x>-1} \Leftrightarrow \heva{&x>-1\\&x\ne 2.}$\\
		Vậy tập xác định của hàm số đã cho là $\mathscr{D}=(-1;+\infty) \setminus \{2\}$.
	}
\end{ex}


\begin{ex}%[Huỳnh Xuân Tín, dự án giảng 10-11, CTGDPT2018]%[0D7C3-6]
	\immini
	{Một chú thỏ ngày nào cũng ra bờ suối ở vị trí $A$ cách cửa hang của mình ở vị trí $B$ là $370$ m để uống nước, sau đó chú thỏ sẽ đến vị trí $C$ cách $A$ $120$ m để ăn cỏ rồi trở về hang. Tuy nhiên, hôm nay sau khi uống nước ở bờ suối, chú thỏ không đến vị trí $C$ như mọi ngày mà chạy đến vị trí $D$ nằm giữa $C$ và $B$ để tìm cà rốt rồi mới trở về hang (xem hình dưới). Biết rằng, tam giác $ABC$ vuông tại $C$, tổng thời gian chú thỏ chạy từ $A$ đến $D$ rồi về hang hết $30$ giây (không kể thời gian tìm cà rốt), tốc độ của chú thỏ trên các đoạn $AD$ và $DB$ lần lượt là $13$ m/s và $15$ m/s. Tính độ dài đoạn $CD$.}
	{\begin{tikzpicture}[line join=round, line cap = round, >=stealth, scale=.8,font=\footnotesize,transform shape]
			\foreach \x/\y/\z/\g in
				{
					0/0/C/-135, 3/0/D/-90, 5/0/B/-45, 0/4/A/135
				}
			\draw[fill=black] (\x,\y) circle(1pt) coordinate (\z) ($(\z)+(\g:3.5mm)$) node{$\z$};
			\draw (D)--(A)--(B)--(C)--(A);
		\end{tikzpicture}}
	\shortans{$50$}
	\loigiai{
		Gọi thời gian chú thỏ chạy trên đoạn $AD$ là $x$ giây ($0<x<30$). Khi đó $AD = 13x$ m.\\
		Thời gian chú thỏ chạy trên đoạn $BD$ là $30-x$ giây nên $BD = 15(30-x)$ m.\\
		Theo định lý Py-ta-go ta có $BC = \sqrt{AB^2 - AC^2} = 350$. \\
		Suy ra $CD = BC - BD = 350 - 15(30-x) = 15x - 100$ m.\\
		Áp dụng định lý Py-ta-go cho tam giác $ACD$ ta có
		\[
			AD^2 = AC^2 + CD^2\Rightarrow (13x)^2 = 120^2 + (15x-100)^2
			\Rightarrow 56x^2 - 3000x + 24400 = 0 \Rightarrow \hoac{&x=10\\ &x = \dfrac{307}{5}.}
		\]
		Mà $0<x<30$ nên $x = 10$, suy ra $CD = 50$.
	}
\end{ex}

\begin{ex}%[VDC8-DuongBuiDuc]
	Gọi $ S $ là tập các số tự nhiên chẵn có $ 4 $ chữ số khác nhau. Lấy ngẫu nhiên một số trong $ S $. Tính xác suất để số được chọn lớn hơn hoặc bằng $ 2018 $.
	\loigiai{
		Gọi số cần tìm có dạng $ \overline{abcd} $ ($ a\neq 0 $).
		\begin{enumerate}[+)]
			\item Ta tìm các số chẵn có $ 4 $ chữ số phân biệt:
			\begin{itemize}
				\item Nếu $ d=0 $ thì số cách sắp xếp cho $ 3 $ vị trí $ a,b,c $ là $ \mathrm{A}^{3}_{9}=504 $ cách.
				\item Nếu $ d\neq 0 $ thì $ d $ có $ 4 $ cách chọn, $ a $ có $ 8 $ cách chọn, $ b $ có $ 8 $ cách chọn, $ c $ có $ 7 $ cách chọn.\\
				Suy ra có $ 4\cdot 8\cdot 8\cdot 7=1792 $ số.
			\end{itemize}
			Vậy tập $ S $ có $ 504+1792=2296 $ số.
			\item Ta tìm các số chẵn có $ 4 $ chữ số phân biệt và nhỏ hơn $ 2018 $:
			\begin{itemize}
				\item Nếu $ \overline{abcd}\geq 2000 $ thì có $ 9 $ số.
				\item Nếu $ \overline{abcd}<2000 $ thì $ a $ có $ 1 $ cách chọn, $ d $ có $ 5 $ cách chọn, $ b $ có $ 8 $ cách chọn, $ c $ có $ 7 $ cách chọn. Suy ra có $ 1\cdot 5\cdot 8\cdot 7=280 $ số.
			\end{itemize}
			Vậy có $ 280+9=289 $ số chẵn có $ 4 $ chữ số phân biệt và nhỏ hơn $ 2018 $, suy ra có $ 2296-289=2007 $ số chẵn có $ 4 $ chữ số phân biệt và lớn hơn hoặc bằng $ 2018 $.
		\end{enumerate}
		Xác suất để chọn được một số lớn hơn hoặc bằng $ 2018 $ từ tập $ S $ là $ \dfrac{2007}{2296} $.
	}
\end{ex}
\Closesolutionfile{ansbook}
